The contribution of an RBP sequence model beyond nucleotide composition depends jointly on the
model, the baseline and the negative-set protocol. In the study panel, the protocol yielding the
highest apparent AUROC left the least incremental signal for the sequence model, while stricter
composition matching lowered apparent AUROC and increased the measured contribution. The
protocol-associated variation was comparable in scale to that produced by changing model
class.

Estimator design also mattered. The conventional two-stage procedure returned a positive
increment for a model whose information was already contained in the baseline. Cross-fitting
the score covariate reduced this null floor by at least 95\% and preserved a 4.84-fold span
across protocols. The externally constructed benchmark reproduced the protocol dependence, but
not the inverse relation between apparent difficulty and contribution.

These findings support a practical reporting standard. Studies should describe the negative-set
construction in reproducible detail, report the composition-only performance under the same
protocol, state the baseline order, and calibrate the incremental estimator with a null model.
Score covariates should be cross-fitted when feasible. Most importantly, contributions measured
under different negative-set protocols should not be interpreted as directly comparable. The
negative set is part of the scientific question posed by the benchmark, not merely a technical
setting used to evaluate it.
